\documentclass[10pt]{article}
\usepackage[margin=1in]{geometry}
\usepackage[T1]{fontenc}
\usepackage[utf8]{inputenc}
\usepackage{lmodern}
\usepackage{microtype}
\usepackage{amsmath,amssymb,amsthm}
\usepackage{graphicx}
\usepackage{booktabs,tabularx,array}
\usepackage{enumitem}
\usepackage[numbers,sort&compress]{natbib}
\usepackage{xcolor}
\usepackage[hidelinks]{hyperref}
\usepackage{url}
\usepackage{placeins}
\newtheorem{proposition}{Proposition}
\newcommand{\Hmethod}{Historical moment}
\newcommand{\Rsmall}{Random R128}
\newcommand{\Rlarge}{Random R138}
\newcommand{\Fmethod}{Fixed width 7}
\newcommand{\analysisdir}{../analysis/confirmation}
\newcommand{\artifact}[1]{\texttt{\detokenize{#1}}}
\setlist{nosep,leftmargin=*}
\setlength{\emergencystretch}{2em}
\title{Do Historical Moments Improve the Risk--Work Frontier?\\
\large A Controlled Study of Scheduled Quadratic Network Growth}
\author{Wenyu Huang\\University of California San Diego}
\date{}

\hypersetup{
  pdftitle={Do Historical Moments Improve the Risk–Work Frontier? A Controlled Study of Scheduled Quadratic Network Growth},
  pdfauthor={Wenyu Huang}
}

\begin{document}
\maketitle
\begin{abstract}
Choosing a promising direction for a new neural unit does not establish an
advantage after finite optimization and charged proposal costs. We study
this distinction in a centered quadratic network whose hidden vectors,
signed output coefficients, and intercept are all trained. A historical
Gaussian moment proposes spectral directions while retaining 128 examples;
the comparators use random growth with 128 or 138 examples, or a fixed
width of seven with 138 examples. Growth from one to seven units follows
a fixed schedule. A locally locked protocol specifies five work budgets,
balanced learning-rate selection on separate development problems, and
256 independent confirmation teachers and streams. Exact population
excess risk is evaluated only after each training trajectory is fixed.
The primary endpoint averages terminal excess risk over a fixed
log-budget grid and requires superiority to all three comparators.
The prespecified conjunction is met ($p=0.00312$): mean integrated
excess risk is $0.00536$, compared with $0.00594$, $0.00579$, and
$0.01498$ for the three controls. The fixed-width mean is driven by
large-error runs; historical moments win on only 118/256 paired seeds.
The effect is modest against random growth, and both random controls already meet the fixed risk target
on every seed at the smallest evaluated budget. Fixed width is faster
in the observed learner timings. Secondary tasks do not show a
uniform historical-moment advantage.
The analytical identities and empirical comparison concern this synthetic
model and optimizer; they supply neither autonomous growth nor a general
risk or runtime guarantee.
\end{abstract}

\section{Question and relation to earlier work}
\label{sec:introduction}

A network can be enlarged without changing its current output. It can
also receive a direction that would reduce risk under an ideal coefficient
refit. The practical question is whether this information improves the
predictor that training actually reaches under a common resource rule.
Direction estimation, replay maintenance, optimizer state, and fitting
all consume resources. An initially better direction can fail to repay
these costs, and a competent fixed-width model may reach the same risk
without any growth mechanism.

This paper studies one narrowly specified version of that question:
does retaining a historical quadratic moment reduce final population
excess risk over a prespecified training-work grid, compared with two
random-growth controls and a fixed-width control? Earlier internal
Phase~5 experiments did not establish that stronger birth directions
produced lower final risk. Phase~6 therefore evaluates a complete budget
frontier using independently rotated teachers, balanced development
selection, and broader paired replication. Its intended contribution
is a controlled empirical finding and an inspectable resource ledger.
The historical-moment mechanism is carried forward as a candidate to
test, not introduced as a new algorithmic principle.

The relevant prior work already separates several ingredients of this
construction. Net2Net develops function-preserving transformations for
transferring a trained network into wider or deeper networks
\cite{net2net}. GradMax chooses initializations for added neurons by
maximizing new-weight gradients and uses a singular-value decomposition
\cite{gradmax}. Growing Tiny Networks formalizes expressivity
bottlenecks relative to functional gradients and uses this information
to add neurons during training \cite{tiny}. In the matrix formulation,
GECO is a greedy method for convex optimization under a low-rank
constraint, using leading singular directions and corrective optimization
\cite{geco}. These are close precedents for function preservation,
gradient-informed expansion, and spectral rank-one search. We do not
claim that these ingredients are new.

The present factorized quadratic model permits exact population
evaluation and transparent small-scale accounting. It is not a
reproduction of those prior algorithms, and the four experimental
comparators do not establish superiority over them. In particular,
finite clipped Adam updates of a nonlinear parameterization differ
from an analytical rank-one coefficient fit or a fully corrective
matrix optimization step. Section~\ref{sec:p6-theory} gives the exact
risk identity and states this separation. The earlier project also
considered a fresh-block matrix recurrence; the discussion there is
included solely to explain why its contraction analysis does not
transfer to this implementation.

% Self-contained section body. The parent document supplies amsmath,
% amssymb, amsthm, and a proposition environment.
% Bibliography keys geco and tiny are already present in phase5 references.
% Primary sources checked: arXiv:1106.1622v1 and arXiv:2405.19816v2.
\section{Mathematical scope and the budget endpoint}
\label{sec:p6-theory}

The mathematical ingredients below are classical Gaussian moment
identities and elementary properties of the chosen evaluation endpoint.
They are not claimed as novel learning principles. Greedy low-rank
optimization and growth based on expressivity bottlenecks already have
established antecedents~\cite{geco,tiny}. Our question is empirical:
does the information supplied by a historical moment improve the
training-budget frontier once its computational costs are charged?

\subsection{Exact population risk of the trained predictor}

Let $X\sim N(0,I_d)$, $H(X)=XX^\top-I_d$, and
$q_A(X)=\langle A,H(X)\rangle$, where matrices are symmetric and
$\langle A,B\rangle=\operatorname{tr}(AB)$. The teacher and predictor are
\begin{equation}
 Y=q_{A_*}(X)+c_*+\xi,
 \qquad f_{B,c}(X)=q_B(X)+c,
 \qquad B=\sum_{j=1}^k a_jw_jw_j^\top,
 \label{eq:p6-model}
\end{equation}
where $\xi$ is independent of $X$, with mean zero and variance
$\sigma^2<\infty$. Our experiments set $c_*=0$ but fit $c$.

\begin{proposition}[Exact risk including the intercept]
\label{prop:p6-risk}
For every symmetric $B$ and scalar $c$,
\begin{equation}
 R(B,c):=\mathbb E[(Y-f_{B,c}(X))^2]
 =\sigma^2+2\|B-A_*\|_F^2+(c-c_*)^2.
 \label{eq:p6-risk}
\end{equation}
Consequently, the excess risk used in the experiments is
$\mathcal E(B,c)=2\|B-A_*\|_F^2+c^2$.
The same formula holds conditionally on a realized training history and
teacher, provided the test example is independent of that history.
\end{proposition}
\begin{proof}
The Gaussian second- and fourth-moment identities give
\[
 \mathbb E[X_iX_j]=\delta_{ij},\qquad
 \mathbb E[X_iX_jX_kX_l]
 =\delta_{ij}\delta_{kl}+\delta_{ik}\delta_{jl}
  +\delta_{il}\delta_{jk}.
\]
Thus $\mathbb E H_{ij}=0$ and
$\mathbb E[H_{ij}H_{kl}]=\delta_{ik}\delta_{jl}
+\delta_{il}\delta_{jk}$. Multiplying by $D_{ij}D_{kl}$ and summing
gives $\mathbb E[q_D(X)^2]=2\|D\|_F^2$ for symmetric $D$.
Take $D=A_*-B$. Since $\mathbb E q_D=0$, the cross term between
$q_D$ and $c_*-c$ vanishes. Noise independence and zero mean eliminate
its cross terms with both quantities. Expanding
$(q_D+c_*-c+\xi)^2$ proves~\eqref{eq:p6-risk}.
Conditioning on a saved predictor and its training history fixes $B,c$;
the independent test input retains the same Gaussian moments.
\end{proof}

The intercept term cannot be omitted: $B=A_*$ and $c=1$, for example,
have excess risk one when $c_*=0$. For each individual run, a separate
evaluator applies \eqref{eq:p6-risk} after its training trajectory has
been fixed. Development and confirmation have distinct roles. Before
observing any new development results, we froze a common search protocol:
all four methods are evaluated at learning rates
$\{0.0075,0.015,0.03\}$ and all five budgets on 16 independent
development teacher/data seeds. For each method, the mean exact
population AUC across these development seeds selects one learning
rate, which is then frozen before the 256-replicate confirmation
experiment. This is explicitly an oracle-based development selection;
its computational cost is reported separately. Confirmation teacher
matrices, population scores and oracle direction coefficients do not
enter the learner's proposal selection, optimization, stopping, or
hyperparameter selection. Exact evaluation removes test-input Monte
Carlo error in this synthetic model; it does not remove variation
across independently generated training problems.

\subsection{Function-preserving birth and the limits of direction quality}

Appending any finite direction $v$ with output coefficient
$a_{k+1}=0$ leaves both $B$ and $c$ unchanged, and hence preserves the
predictor for every input. Initializing the new Adam moment states to
zero also leaves the function unchanged. The subsequent optimization
trajectory is a separate matter: at birth the gradient with respect to
the new hidden vector is zero because it is multiplied by its output
coefficient, while the new output coefficient can have a nonzero
gradient.

For a unit direction $v$, set $E=A_*-B$ and hold all other parameters,
including the intercept, fixed. Proposition~\ref{prop:p6-risk} yields
\begin{equation}
 R(B+\alpha vv^\top,c)-R(B,c)
 =2\alpha^2-4\alpha v^\top Ev.
 \label{eq:p6-direction-gain}
\end{equation}
The oracle scalar coefficient is $\alpha=v^\top Ev$, with improvement
$2(v^\top Ev)^2$. Maximizing this quantity over unit directions gives
$2\|E\|_{\mathrm{op}}^2$. This is the improvement available from an
ideal scalar refit at a fixed predictor, not an improvement guaranteed
by the implemented factor Adam updates.

In particular, the earlier project's fresh-block quadratic analysis
concerns a prescribed update of the form
\[
 B^+=B+\frac{1}{2n}\sum_{i=1}^n
       \{Y_i-q_B(X_i)\}H(X_i),
\]
under its stated zero-intercept model and conditional independence
assumptions. Factor Adam instead updates $(w_j,a_j,c)$ using empirical
gradients, accumulated coordinatewise moments, and clipping. It may
reuse retained examples and historical summaries. Neither its update
map nor its sample dependence matches that matrix recurrence, so the
contraction theorem does not supply a convergence rate or cost bound
for the present experiment. Even the exact empirical residual matrix,
when the predictor has an intercept, is
\begin{align*}
 G_n(B,c)&=S_n-T_n(B)-\tfrac12c\overline H_n,
 &S_n&=\frac1{2n}\sum_iY_iH_i,\\
 T_n(B)&=\frac1{2n}\sum_iq_B(X_i)H_i,
 &\overline H_n&=\frac1n\sum_iH_i.
\end{align*}
The historical proxy $S_n-B$ uses the known population Gaussian
operator; it is not this exact empirical residual.

\subsection{A fixed log-budget area from five completed runs}

The five preregistered per-arrival budget caps, in the declared work
units, are
\[
 (b_1,b_2,b_3,b_4,b_5)
 =10^6(3.2,5,6.4,8,10).
\]
Every run processes seven arrivals under the same data-availability
rule. Let $e_{mrj}$ be the final excess risk of method $m$, replicate
$r$, and budget setting $b_j$, after all seven arrivals. Each entry is
obtained from its own completed training run. These five entries are
not checkpoints extracted from one run with the largest budget.

Write $x_j=\log b_j$ and $L=x_5-x_1>0$. Define the normalized
log-budget trapezoid area
\begin{equation}
 \operatorname{AUC}_{mr}
 =\frac1L\sum_{j=1}^{4}(x_{j+1}-x_j)
                   \frac{e_{mrj}+e_{mr,j+1}}2
 =\sum_{j=1}^5\omega_j e_{mrj},
 \label{eq:p6-auc}
\end{equation}
where
\[
 \omega_1=\frac{x_2-x_1}{2L},\qquad
 \omega_5=\frac{x_5-x_4}{2L},\qquad
 \omega_j=\frac{x_{j+1}-x_{j-1}}{2L}\quad(2\le j\le4).
\]
The weights are positive, sum to one, and are fixed independently of
all outcomes. Only budget is logarithmically transformed; the ordinate
is excess risk itself. Lower area is better. Replacing $b_j$ by its
seven-arrival cap $7b_j$ leaves all weights unchanged, as does changing
the logarithm base.

Equation~\eqref{eq:p6-auc} is an integral of the piecewise-linear
interpolant in log budget, not an assumption that the unobserved
performance curve is linear. The budget axis is the assigned cap;
actual charged work, unused terminal budget and wall time are reported
separately. The area therefore summarizes the specified allocation
protocol and cost ledger, and does not by itself establish hardware
runtime dominance. The primary sample comprises 256 independent
teacher/data replicates, with pairing across methods and budget
settings within each replicate. The five budgets do not multiply the
number of independent inference units by five.

\subsection{Meaning of the three-comparison intersection--union test}

Let $h$ denote historical-moment growth, and let $m_1,m_2,m_3$ denote
the three preregistered primary comparators. For $\ell=1,2,3$, define
the paired area difference and its mean over the specified distribution
of independently generated problems by
\[
 D_{r\ell}=\operatorname{AUC}_{hr}
                -\operatorname{AUC}_{m_\ell r},
 \qquad \mu_\ell=\mathbb E[D_{r\ell}].
\]
The three-comparison intersection--union test (IUT3) addresses
\begin{equation}
 H_0:\ \bigcup_{\ell=1}^3\{\mu_\ell\ge0\},
 \qquad
 H_1:\ \bigcap_{\ell=1}^3\{\mu_\ell<0\}.
 \label{eq:p6-iut}
\end{equation}
If $p_\ell$ is a valid one-sided $p$-value for the corresponding
component null, reject~\eqref{eq:p6-iut} at level $\alpha$ only when
all three $p_\ell\le\alpha$, equivalently when
$p_{\rm IUT}=\max_\ell p_\ell\le\alpha$.
Indeed, under the union null at least one component null, say
$\ell_0$, is true, and
\[
 \mathbb P(p_{\rm IUT}\le\alpha)
 \le\mathbb P(p_{\ell_0}\le\alpha)\le\alpha.
\]
This argument requires no independence between the three comparisons
and no Bonferroni division of $\alpha$ for this single conjunction
claim. Its calibration still depends on the validity of each chosen
component test; a finite-sample exact claim cannot be inferred merely
from having 256 replicates. Additional families of claims or
outcome-dependent selection of comparators require their own handling.

Rejection supports lower mean area against every one of the three
specified comparators under this protocol. It does not imply lower
risk at every budget, success on every seed, superiority to untested
methods, or optimality among architectures. Failure to reject does
not establish equivalence, absence of a directional effect, or absence
of an advantage at some budgets. A superiority test has no prespecified
equivalence margin. Together, the risk identity, direction diagnostic
and IUT3 interpretation provide no uniform optimal cost bound for
factor Adam or autonomous network growth.


\section{Experimental protocol}
\label{sec:methods}

\subsection{Data stream and learner information}

Each primary replicate uses $d=12$ and seven successive arrivals of
1,024 independent observations, for 7,168 observations in total.
The teacher is $A_*=Q\Lambda Q^\top$, where $Q$ is obtained from an
independent Gaussian QR draw and
\[
\operatorname{diag}(\Lambda)
=(0.65,-0.50,0.30,-0.20,0,\ldots,0).
\]
Inputs are independent $N(0,I_{12})$, the teacher intercept is zero,
and additive Gaussian noise has standard deviation $0.3$. Each
replicate has its own teacher rotation, observations, noise, and
optimization random streams. Methods and budgets are paired within
the same replicate. At each arrival, a method accesses only the current
block and its permitted retained state; it never accesses future
observations and never receives the teacher matrix or a population
score while training.

At an arrival, the fitting pool consists of the fresh 1,024 observations
and the reservoir retained from previous arrivals. A reservoir update
then processes the fresh block using Algorithm~R. The fresh observations
are not duplicated in that arrival's fitting pool merely because some
of them enter the updated reservoir. Independent random namespaces
separate data generation, initialization, proposal directions, fitting
permutations, reservoir selection, and execution order. Within a seed,
all trajectories are completed before a separate evaluator adds their
exact population risks. The evaluator's access to $A_*$ is an
evaluation permission, not a learner input.

\subsection{Four complete algorithms}

The prediction function implemented in all four methods is
\begin{equation}
 f_\theta(x)=c+\sum_{j=1}^k a_j\bigl((w_j^\top x)^2-\|w_j\|_2^2\bigr).
 \label{eq:factor-predictor}
\end{equation}
All $w_j$, $a_j$, and $c$ are trainable; the comparison is not a
fixed-feature regression. At initialization, hidden vectors are
independent Gaussian draws normalized to unit length, output
coefficients are drawn from $N(0,0.01^2)$, and $c=0$.
The three growing methods start at width one. Before fitting arrivals
two through seven, each appends exactly one hidden unit. Its output
coefficient and Adam moments are zero, while existing parameters and
states are preserved. The fixed-width method initializes seven units
and optimizes all of them from the first arrival.

\Hmethod\ retains a full $12\times12$ matrix sum and a reservoir of
128 examples. After incorporating the fresh block, let $N_t$ denote
the number of observations received and define
\begin{equation}
 S_t=\frac{1}{2N_t}\sum_{i\le N_t}Y_i(X_iX_i^\top-I_d),\qquad
 M_t=S_t-B_{t^-}.
 \label{eq:moment-proposal}
\end{equation}
The stored sum is updated once per observation; no historical raw
data beyond the reservoir are retained by the learner. The new
direction is a unit eigenvector of $M_t$ associated with its largest
absolute eigenvalue. Both signs of teacher eigenvalues matter because
the output coefficients are signed. The Gaussian population identity
$\mathbb E S_t=A_*$ motivates this proxy, but $M_t$ is not the exact
empirical residual matrix of the fitted predictor with intercept.

\Rsmall\ and \Rlarge\ use normalized independent Gaussian birth
directions and retain 128 and 138 examples, respectively. The former
controls for the number of replay examples; the latter spends
approximately the historical matrix's storage on additional examples.
\Fmethod\ retains 138 examples. No method uses a gate, a validation
controller, an accepted/rejected proposal branch, or an autonomous
choice of birth timing or width.

\begin{table}[t]
\centering\small
\begin{tabular}{lclrr}
\toprule
Method & Width & Birth direction & Replay rows & Retained bytes\\
\midrule
\Hmethod & $1\to7$ & Historical spectral & 128 & 15,488\\
\Rsmall & $1\to7$ & Random & 128 & 14,336\\
\Rlarge & $1\to7$ & Random & 138 & 15,456\\
\Fmethod & 7 & No births & 138 & 15,456\\
\bottomrule
\end{tabular}
\caption{Persistent observation-dependent arrays after filling the
reservoir. Every stored row has 12 float64 inputs, a float64 label,
and an int64 observation identifier. The historical matrix adds
144 float64 words. Its 32-byte difference from the 138-row controls
is disclosed; storage is close but not exactly matched. Model,
optimizer, working arrays, and control state are separate.}
\label{tab:methods}
\end{table}

\subsection{Fitting and optimizer-age sensitivity}

The empirical squared-error gradients update all parameters using
Adam with $(\beta_1,\beta_2)=(0.9,0.999)$, numerical stabilizer
$10^{-8}$, and global gradient-norm clipping at 20. The batch size is
128. Random permutations are traversed with tails carried into the
next permutation, so an epoch boundary alone does not create a short
optimizer batch. If the remaining budget cannot fund 128 examples,
the largest affordable positive final batch is used and recorded.
Training stops when no positive batch fits the remaining declared
work. There is no validation-based early stopping or rollback.

The primary optimizer uses the global Adam update count for bias
correction, with zero moment states for newborn coordinates. This
choice gives new and old coordinates different histories despite a
shared global count. A separately specified sensitivity uses each
coordinate's local age for bias correction; newborn ages start at
zero and existing ages continue. The sensitivity repeats the first
32 primary seeds with frozen learning rates. Its results are paired
with those same seeds' primary runs and are not 32 additional
independent primary replicates.

\subsection{Allocated budgets and charged operations}
\label{sec:accounting}

Each method runs independently at each of the five per-arrival caps
$10^6(3.2,5,6.4,8,10)$. Thus the full-stream allocated totals are
$10^6(22.4,35,44.8,56,70)$. A higher-budget run does not continue a
lower-budget run. Before fitting each arrival, the implementation
reserves declared costs for initialization when applicable,
observation identifiers, retention allocation and maintenance,
fitting-pool copies, historical-moment maintenance, residual-matrix
construction, eigendecomposition, direction construction, births,
and event control. The remaining allocation funds fitting and
permutation operations.

Work is a deterministic scalar-operation proxy defined by explicit
formulas in the runner. It is not measured hardware FLOPs.
For a batch of $n$ examples at width $k$, let $P=(d+1)k+1$.
The fitting charge is
\begin{equation}
 C_{\mathrm{fit}}(n,d,k)
 =n(6dk+12k+2d+12)+30P+4dk+4k+32.
 \label{eq:work-fit}
\end{equation}
Each newly generated permutation of a fitting pool of size $L$
costs $4L$. A moment update costs
$2nd^2+nd+n+4d^2$, and an eigendecomposition costs the declared
surrogate $10d^3$. Additional formulas and the separation of
measurement categories are given in Appendix~\ref{sec:repro}.
All methods pay their own charges under the same definitions, so
equal allocated work need not yield equal optimizer steps. Every
arrival records actual charged work and its integer remainder.

Learner timings sum algorithm blocks covering initialization, proposal,
retention, optimization, and cleanup. Data generation, post-trajectory
evaluation, archive serialization, and major audit-counter updates
are excluded from those blocks; minor control, allocation, and
instrumentation overhead remains. The code inventories retained arrays,
parameter and optimizer arrays, and estimated explicit working
arrays. These are neither operating-system peak resident memory nor
a measurement of NumPy/BLAS internal allocations. Wall times depend
on the implementation, process contention, memory allocation, and
the host; a declared-work advantage is not automatically a
wall-clock advantage.

\subsection{Separate development selection and primary inference}

Before confirmation, every method receives the same selection rule:
16 development seeds, all five budgets, and three learning rates,
$\{0.0075,0.015,0.03\}$. The method-specific rate minimizing mean
development AUC is selected, with the smaller rate breaking an exact
tie. This entails 960 complete development trajectories. Development
uses independently generated teachers and streams, and its exact
population scores are an explicitly disclosed selection oracle.
The selected rates are frozen for confirmation, secondary tasks,
and the local-age sensitivity. Development work is a one-time
experimental selection cost, disclosed separately from a single
trained instance's work.

The primary confirmation sample comprises exactly 256 seeds. For each
comparator we compute the seedwise difference in the AUC defined in
Equation~\eqref{eq:p6-auc}. The primary decision requires all three
one-sided paired-$t$ upper 95\% confidence bounds to lie below zero.
This is the intersection--union test of
Equation~\eqref{eq:p6-iut}. We report its maximum component $p$-value,
each paired standard deviation and interval, and 10,000-resample
paired bootstrap intervals as a sensitivity check. The $t$
approximation is not a finite-sample guarantee for unbounded
squared loss. A conservative planning calculation at $N=256$
requires roughly a $0.20$ standardized effect in the weakest
contrast for 80\% joint power. This is a planning approximation,
not retrospective evidence of power or equivalence.

The protocol was locally locked before new smoke or development
outcomes, with a subsequent code freeze following correctness review.
It was not externally registered or timestamp-certified. The seeds,
budgets, primary statistic, learning-rate candidates, and sample size
were fixed prospectively. Smoke runs serve correctness and timing
checks only. Necessary implementation corrections and their source
snapshots are preserved. Poor-performing seeds remain in the sample;
nonfinite or incomplete runs require an explicit investigation
instead of silent exclusion. Confirmation results do not authorize
an additional seed, budget, task, or optimizer search.

\subsection{Secondary questions}

The fixed target for the primary noisy task is excess risk
$\epsilon=0.01$. We report attainment at every budget and the first
successful evaluated final-budget point, retaining failures as
nonattainment. This is a discrete five-point grid summary. It does
not assume a monotone risk curve, interpolate an unmeasured crossing,
or estimate the earliest training-time crossing.

Two predeclared compression comparisons place the historical method
at 80\% of the comparator's allocated work: per-arrival 6.4 versus
8 million units and 8 versus 10 million units. We report paired
terminal-risk differences and target attainment. These exploratory
comparisons do not establish a unique 20\% savings, because a
comparator might also reach the target at a lower budget.

Two descriptive controls each use 32 independent seeds and all four
methods at all five budgets. The null teacher has $A_*=0$ and noise
standard deviation $0.5$; it probes scheduled expansion when no
quadratic signal is present. The noiseless rank-two teacher has
eigenvalues $0.8$ and $0.45$ with independent rotation; it probes
dependence on label noise and spectrum. Neither control supplies
additional primary-test replications.

\section{Results}
\label{sec:results}

\subsection{Frozen selection and primary frontier}
\input{\analysisdir/results_text.tex}
\input{\analysisdir/primary_table.tex}
\IfFileExists{\analysisdir/figures/indefinite_noisy_global_frontier.pdf}{
\begin{figure}[t]
\centering
\includegraphics[width=\linewidth]{\analysisdir/figures/indefinite_noisy_global_frontier.pdf}
\caption{Primary final-risk frontier and seed-paired differences after
seven arrivals. Each budget setting is trained independently for every
seed; plotted points are means over 256 seeds. Lines connect evaluated
budget settings. Error bars are descriptive
95\% intervals across the 256 paired problems, not simultaneous
confidence bands. Only the allocated-work axis is logarithmic.}
\label{fig:primary-frontier}
\end{figure}}{}

Mean primary AUC is $0.0053573$ for historical moments, $0.0059388$
for R128, $0.0057884$ for R138, and $0.014979$ for fixed width seven.
The historical reductions relative to the two random controls are
about 9.8\% and 7.4\% of their \emph{excess-risk areas}. They are not
reductions of total risk, allocated work, or elapsed time. The paired
standard deviations are $0.0006434$, $0.0009075$, and $0.05582$,
respectively. The fixed-width contrast is therefore much more variable
and tail-sensitive than the random-growth contrasts. A descriptive
audit of its distribution shows median AUC $0.005112$, slightly below
the historical median $0.005237$, and maximum AUC $0.60066$.
Historical moments have lower paired AUC on 118/256 problems and
higher AUC on 138; the median paired difference is $+0.00008489$.
The five problems most favorable to historical moments account for
68.34\% of the net mean difference. Thus this contrast concerns
\emph{expected mean risk and large-error outcomes}, not a typical-seed
or majority-seed advantage. These post-analysis distribution
diagnostics do not alter the primary endpoint or exclude any seed.
The paired bootstrap interval remains negative, but it resamples
only the observed tail. Neither that interval nor the $t$
approximation supplies distribution-free finite-sample coverage or
resolves uncertainty about unobserved large errors.

The average-area result does not imply pointwise dominance. At the
smallest cap, historical moments have mean excess risk $0.005514$,
versus $0.005284$ and $0.005126$ for R128 and R138. The historical
method has a lower mean than both random controls at the other four
caps. Its mean risk is lowest at the 5-million per-arrival cap and
then rises, reaching $0.005940$ at 10 million. The observed curves
are not monotone training-efficiency frontiers. They describe terminal
performance under fixed-rate fitting and a fixed arrival schedule;
there is no selected best-checkpoint envelope.

Figure~\ref{fig:distribution-diagnostic} makes the fixed-width upper
tail and the smallest-budget reversal against random growth visible
using all 256 primary seeds. This descriptive post-analysis
visualization adds no hypothesis test and excludes no observation.

\begin{figure}[t]
\centering
\includegraphics[width=\linewidth]{\analysisdir/figures/primary_distribution_diagnostic.pdf}
\caption{Descriptive post-analysis visualization of all 256 primary
seeds. Left: empirical cumulative distributions of seed-level AUC,
including every large fixed-width value. The logarithmic display
axis does not change the raw-excess-risk primary endpoint. Right:
historical minus random-growth terminal-risk differences at each
allocated budget, with ordinary paired 95\% $t$ intervals. Positive
differences at the smallest budget favor the random controls; these
pointwise intervals are not simultaneous confidence bands.}
\label{fig:distribution-diagnostic}
\end{figure}

All methods selected $0.0075$, the smallest candidate learning rate.
Thus balanced selection was completed as planned, but the development
grid does not establish globally optimized controls or rule out
improvement from lower rates. The locked confirmation was completed
without extending that grid after observing the result.

\subsection{Target attainment, charged resources, and sensitivity}
At the smallest evaluated cap, the $0.01$ target is attained by
255/256 historical runs, 256/256 R128 runs, 256/256 R138 runs, and
237/256 fixed-width runs. Across the complete evaluated grid, the
corresponding counts are 256, 256, 256, and 248; eight fixed-width
problems never attain it. First attainment is already left-censored
by the grid for nearly all growing runs, so the target cannot identify
a historical-moment work saving. After their first successful grid
point, 3 historical, 16 R128, and 19 R138 problems fail at a later
point. This supports reporting every evaluated outcome instead of
assuming monotonicity.

The 6.4-versus-8-million compression comparison favors historical
moments in mean excess risk against R128, R138, and fixed width:
the differences are $-0.000791$, $-0.000747$, and $-0.008663$,
with descriptive 95\% paired intervals
$[-0.000948,-0.000634]$, $[-0.000929,-0.000566]$, and
$[-0.015640,-0.001687]$. In the 8-versus-10-million comparison,
the differences are $-0.000740$, $-0.000539$, and $-0.004498$;
the fixed-width interval $[-0.009526,0.000529]$ crosses zero.
These comparisons use approximately 0.8000 times the charged work,
but do not demonstrate 20\% minimum-cost savings. In particular,
both random controls already reach the target at the lowest cap.

\begin{table}[t]
\centering\small
\begin{tabular}{lrrrr}
\toprule
Task / age ($N$) & Historical & R128 & R138 & Fixed 7\\
\midrule
Primary / global (256) & 0.005357 & 0.005939 & 0.005788 & 0.014979\\
Null / global (32) & 0.008164 & 0.009189 & 0.008923 & 0.007243\\
Rank two / global (32) & $1.142\!\times\!10^{-6}$ & $7.639\!\times\!10^{-7}$ & $8.701\!\times\!10^{-7}$ & $1.091\!\times\!10^{-4}$\\
Primary / local (32) & 0.004689 & 0.005483 & 0.005675 & 0.024995\\
\bottomrule
\end{tabular}
\caption{Mean normalized terminal excess-risk areas. Only the first row
belongs to the primary inference. The local-age row reuses the first
32 primary problems; it is not a comparison with all 256 global-age
runs. Full intervals and paired contrasts are retained in the
analysis CSVs and generated detailed table.}
\label{tab:groupmeans}
\end{table}

On the same 32 primary problems, replacing global by local ages
reduces the historical method's mean AUC by $0.0006833$
(descriptive 95\% paired interval $[-0.0008481,-0.0005184]$).
Fixed-width trajectories are unchanged. Under local age, historical
minus R128, R138, and fixed-width mean differences remain negative,
at $-0.0007941$, $-0.0009863$, and $-0.02031$; their ordinary
two-sided $t$ intervals cross zero for R138 and fixed width. The
direction of the means survives this sensitivity, but its smaller
sample does not establish the same three-way precision. It is not
used to replace or improve the prespecified primary decision.

The task controls show the scope of the result. With the null teacher,
historical minus fixed-width AUC is $0.0009212$
($[-0.0001335,0.001976]$). With the noiseless rank-two teacher,
historical AUC exceeds R128 by $3.783\times10^{-7}$
($[1.569\times10^{-8},7.409\times10^{-7}]$), while all growing
methods reach very small absolute errors. Thus the positive primary
result is not a uniform ranking across even these two nearby
synthetic tasks. These intervals are descriptive.

\begin{table}[t]
\centering\small
\begin{tabular}{lrrr}
\toprule
Method & Mean charged work ($10^6$) & Mean Adam steps & Mean seconds\\
\midrule
Historical moment & 45.63859 & 1,144.0 & 0.08494\\
Random R128 & 45.63651 & 1,203.2 & 0.08840\\
Random R138 & 45.63806 & 1,203.4 & 0.08840\\
Fixed width 7 & 45.63838 & 547.4 & 0.04790\\
\bottomrule
\end{tabular}
\caption{Primary-run resources, averaged equally over 256 seeds and
five independent budget settings per method. Mean allocated work
is 45.64 million units. Seconds sum timed learner blocks for one
complete trajectory; they are implementation observations under
concurrent execution, not a controlled hardware benchmark.}
\label{tab:resources}
\end{table}

The primary study charges $2.33666\times10^{11}$ work units across
5,120 trajectories, against $2.336768\times10^{11}$ allocated.
Total unused work per trajectory ranges from 413 to 7,269 units;
every arrival respects its cap. The historical method performs fewer
Adam steps than random growth because its summary and proposal
charges leave less work for fitting. The comparison therefore
concerns complete algorithms, not an isolated intervention on the
birth direction. Summed primary learner time is 396.35 seconds,
while elapsed invocation time is 215.19 seconds under concurrent
execution. Fixed width is faster in these observed timings despite
its higher primary mean risk, so the result does not establish
wall-clock dominance. Secondary and age-sensitivity invocations add
54.55 and 29.55 elapsed seconds, respectively, with their costs
retained separately.

\section{Discussion and limits}
\label{sec:discussion}

The locked study detects a narrow positive finding: historical-moment
growth has lower mean terminal excess-risk area than all three
specified controls on the indefinite noisy task. Its modest advantage
over random growth remains after charging moment maintenance and
search. The local-age means point in the same direction, with wider
uncertainty. The work-target evidence and wall times do not establish
a minimum-cost or general runtime advantage. These qualifications
are part of the finding, rather than secondary exceptions to a broad
efficiency claim.

The comparison with fixed width is specifically an expected-mean
statement: fixed width has the lower paired AUC on a majority of
primary seeds, while a few severe fixed-width errors dominate the
mean difference. A reader valuing typical-case performance may
therefore draw a different practical conclusion from the same
complete records.

The experimental question concerns final predictors after all seven
arrivals. A favorable newborn direction, a function-preserving birth,
and a final-risk improvement are different statements. The first can
be measured at a particular state; the second is an algebraic
property; the third depends on the subsequent optimization and
resource allocation. An analytical best-coefficient gain cannot
replace the third measurement.

The two random controls also answer different questions. R128 holds
the replay capacity fixed while changing direction information and
the work needed to obtain it. R138 instead provides additional raw
examples for approximately the same persistent data-summary storage.
Neither comparison alone isolates every mechanism. Fixed width seven
tests whether this scheduled growth procedure is useful relative to
a full-width model trained from the start with a separately selected
learning rate. Because earlier growing stages are narrower, their
per-batch charges and optimizer update counts differ from fixed width.
Those differences are part of the specified algorithms, rather than
an unmeasured free benefit.

The primary endpoint is an average over a selected allocation grid.
The positive finding supports lower mean integrated terminal
excess risk against the three tested comparators under this protocol.
It does not establish pointwise dominance, generalization
to other spectra or input distributions, or a hardware-independent
efficiency theorem. The local-age control helps expose sensitivity
to the treatment of newborn coordinates, but it is smaller and uses
rates selected under the global-age primary optimizer.

The known Gaussian input law is a substantial modeling advantage:
it gives both the population operator behind the proposal and the
exact risk evaluator. Non-Gaussian inputs, distribution drift,
unmodeled nonlinear targets, and large-dimensional networks may
change both. Memory scales quadratically with dimension for the
historical matrix, whereas a retained row scales linearly. At $d=12$
the full eigendecomposition is inexpensive enough to study locally;
the declared $10d^3$ charge is not evidence about a large-scale
spectral implementation. The three real neural-growth papers and
GECO cited above were not reimplemented as resource-matched
baselines. A general comparative claim would require that additional
work and an appropriate benchmark.

The bounded experiment is complete. Although the primary conjunction
passes, the left-censored target comparison, faster fixed-width
timings, nonmonotone budget curves, and task dependence do not justify
promoting this finding to a broad efficiency claim. A detectable AUC
difference alone does not establish a new learning principle. No
additional search was undertaken to make the secondary outcomes
more favorable, and intervals containing zero remain inconclusive
for superiority rather than evidence of equivalence.

\section{Reproducibility and artifact scope}

The accompanying local artifact contains the prospectively specified
protocol, frozen code, configuration and seed lists, generated teacher
and observation arrays, per-arrival parameters and optimizer states,
proposal records, data-access ledgers, raw results, independent audit
reports, analysis tables, and vector figures. The paper's numerical
tables are generated from saved results rather than transcribed from
plots. Code and configuration hashes connect each run to its
implementation. The host is an Apple M1 Pro with eight logical CPUs
and 16~GiB of memory; experiments use the existing NumPy CPU
environment with BLAS thread counts pinned to one. Reproduction
commands and exact environment metadata are provided with the
artifact. The files constitute a local research record, not an
external certification of execution or an already published study.

Independent saved-artifact checks passed for all 7,040 primary,
secondary, and age-sensitivity trajectories: 49,280 arrival records
and 112,640 saved parameter-risk evaluations were inspected, with
maximum risk-reconstruction discrepancy $1.11\times10^{-15}$.
Separate checks cover analytical gradients, exact Gaussian quadrature,
future-data isolation, unchanged non-timing trajectories under
execution-order reversal, budget boundaries, optimizer ages, and
deliberate artifact-corruption fixtures. Audits establish consistency
of these saved records and implementation checks; they do not
validate an untested statistical generalization. Earlier project
files were protected by a 1,936-file SHA-256 baseline and remained
unchanged.

\bibliographystyle{plainnat}
\bibliography{references}

\appendix
\section{Resource formulas and audit map}
\label{sec:repro}

Table~\ref{tab:costs} gives the nonfitting charges used with
Equation~\eqref{eq:work-fit}. Here $n=1{,}024$, $r$ is the previously
retained count, $C$ is reservoir capacity, and $k^-,k$ are widths
before and after birth. Each arrival's charges respect its cap;
unused remainders are recorded separately.

\begin{table}[ht]
\centering\small
\begin{tabularx}{\linewidth}{lX}
\toprule
Operation & Declared scalar-work charge\\
\midrule
Moment update & $2nd^2+nd+n+4d^2$\\
Historical residual matrix & $dk^-+d^2(2k^- -1)+2d^2$\\
Eigendecomposition & $10d^3$\\
Historical / random direction & $4d$ / $8d$\\
Reservoir maintenance & $n(d+10)$\\
Fitting-pool copy & $(r+n)(d+2)$\\
Arrival identifiers & $n$\\
Model initialization & $8dk+k+5P$\\
Retention allocation & $C(d+2)$, plus $d^2$ for a historical matrix\\
Birth and state copies & $4P+5(d+1)$\\
Event control & 64\\
New fitting permutation & $4(r+n)$\\
\bottomrule
\end{tabularx}
\caption{Accounting formulas charged when the corresponding operation
occurs. Initialization and retention allocation occur only once;
birth-related charges occur only at scheduled births. Reservoir
maintenance conservatively charges each incoming observation,
including those not retained.}
\label{tab:costs}
\end{table}

The $P=(d+1)k+1$ parameters, Adam first and second moments, and
coordinate ages occupy $32P$ bytes: 92 parameters and 2,944 bytes
at width seven. Retained-data bytes appear in Table~\ref{tab:methods}.
Working arrays are estimated separately. Python and control state,
vendor-internal allocations, data-owner arrays, and audit archives
are excluded; this inventory does not establish peak-memory matching.

\end{document}
